% ---
% id: protocol-compiler/req:feasibility-witness
% kind: req
% title: "Default candidate before a supplied witness"
% summary: >-
%   At a non-empty constraint list the engine tests the default candidate
%   first and, if it fails, a supplied witness, emits neither, and refuses
%   exactly when the default candidate fails and no supplied witness verifies.
% prominence: standard
% status: implemented
% uses: [thm:feasibility, def:validator]
% mentions: [def:feasible]
% assumes: []
% refs: []
% provenance: newly-derived
% implements: certiqs_key.protocol_compiler.feasibility:check_feasibility
% lean:
% ---

Theorem~\ref{thm:feasibility}\,(ii) reduces D(vii) at a non-empty constraint list to
exact decisions on at most two candidates, and check V7 of
Definition~\ref{def:validator} is where they are taken.

\begin{requirement}\label{req:feasibility-witness}
Where the constraint list is non-empty, on a candidate declaration
satisfying D(i)--D(iii) and D(vi), which checks V1--V3 and V6 of
Definition~\ref{def:validator} decide (a candidate failing one of them is
refused there), the engine decides the
conditions of
\eqref{eq:feas}, in the general format those of \eqref{eq:feas-general},
first on the default candidate \eqref{eq:witness-empty}, in the general format
\eqref{eq:witness-general}, and, if it fails, on a
supplied witness $\omega^{\star}$, read from the input
\iface{feasibility_witness}. It forms the default candidate only to test it,
and it emits neither candidate. It refuses with \iface{no_feasibility_witness}
exactly when the default candidate fails and either no witness is supplied or
the supplied one fails.
\end{requirement}
